% Options for packages loaded elsewhere
\PassOptionsToPackage{unicode}{hyperref}
\PassOptionsToPackage{hyphens}{url}
\documentclass[
  12pt,
]{article}
\usepackage{xcolor}
\usepackage{amsmath,amssymb}
\setcounter{secnumdepth}{-\maxdimen} % remove section numbering
\usepackage{iftex}
\ifPDFTeX
  \usepackage[T1]{fontenc}
  \usepackage[utf8]{inputenc}
  \usepackage{textcomp} % provide euro and other symbols
\else % if luatex or xetex
  \usepackage{unicode-math} % this also loads fontspec
  \defaultfontfeatures{Scale=MatchLowercase}
  \defaultfontfeatures[\rmfamily]{Ligatures=TeX,Scale=1}
\fi
\usepackage{lmodern}
\ifPDFTeX\else
  % xetex/luatex font selection
\fi
% Use upquote if available, for straight quotes in verbatim environments
\IfFileExists{upquote.sty}{\usepackage{upquote}}{}
\IfFileExists{microtype.sty}{% use microtype if available
  \usepackage[]{microtype}
  \UseMicrotypeSet[protrusion]{basicmath} % disable protrusion for tt fonts
}{}
\makeatletter
\@ifundefined{KOMAClassName}{% if non-KOMA class
  \IfFileExists{parskip.sty}{%
    \usepackage{parskip}
  }{% else
    \setlength{\parindent}{0pt}
    \setlength{\parskip}{6pt plus 2pt minus 1pt}}
}{% if KOMA class
  \KOMAoptions{parskip=half}}
\makeatother
\usepackage{longtable,booktabs,array}
\usepackage{calc} % for calculating minipage widths
% Correct order of tables after \paragraph or \subparagraph
\usepackage{etoolbox}
\makeatletter
\patchcmd\longtable{\par}{\if@noskipsec\mbox{}\fi\par}{}{}
\makeatother
% Allow footnotes in longtable head/foot
\IfFileExists{footnotehyper.sty}{\usepackage{footnotehyper}}{\usepackage{footnote}}
\makesavenoteenv{longtable}
\setlength{\emergencystretch}{3em} % prevent overfull lines
\providecommand{\tightlist}{%
  \setlength{\itemsep}{0pt}\setlength{\parskip}{0pt}}
\usepackage{bookmark}
\IfFileExists{xurl.sty}{\usepackage{xurl}}{} % add URL line breaks if available
\urlstyle{same}
\hypersetup{
  hidelinks,
  pdfcreator={LaTeX via pandoc}}

\author{SCX}
\date{2026}

\begin{document}

\section{Situs (Physical Positional Encoding)
Physical Realism Validation and Applicability Scenario Analysis}\label{situs-physical-positional-encoding-ux7269ux7406ux771fux5b9eux6027ux9a8cux8bc1ux4e0eux9002ux7528ux573aux666fux5206ux6790}

\begin{center}\rule{0.5\linewidth}{0.5pt}\end{center}

\textbf{Document Status}: Honest Critique --- No Euphemisms\\
\textbf{Prerequisite Reading}:
\texttt{ppe\_rigorous\_derivation.md}（Rigorous Mathematical Derivation),\texttt{multi\_head\_spring\_and\_positional\_encoding\_analysis.md}（CC
Audit Report)\\
\textbf{Methodology}: Real Physical Data + Physical Correspondence of Mathematical Theorems + Scenario-by-Scenario Argumentation\\
\textbf{Core Criterion}: 定理 2.2.1 ---
\(\delta_s^{\text{PE}} > 0 \iff I(Y; P \mid S) > 0\)，即PhysicsLocation在给定Status原子Conditions下必须携带关于标签的额外信息

\begin{center}\rule{0.5\linewidth}{0.5pt}\end{center}

\section{Part 1
Part: Physical Realism Check}\label{ux7b2c-1-ux90e8ux5206ux7269ux7406ux771fux5b9eux6027ux68c0ux67e5}

\begin{center}\rule{0.5\linewidth}{0.5pt}\end{center}

\subsection{1.1
蛋白质序列Location编码}\label{ux86cbux767dux8d28ux5e8fux5217ux4f4dux7f6eux7f16ux7801}

\subsubsection{1.1.1 相同氨基酸，不同Location：Lys-37（Active Site）vs
Lys-289（Surface
loop）}\label{ux76f8ux540cux6c28ux57faux9178ux4e0dux540cux4f4dux7f6elys-37ux6d3bux6027ux4f4dux70b9vs-lys-289ux8868ux9762-loop}

\textbf{Biological Facts (Indisputable)}：

同一个赖氨酸残基在蛋白质不同Location具有完全不同的功能Role： -
\textbf{Active Site Lys-37}：典型催化残基，侧链 pKa 可偏离溶液值 2--4
个单位（From \textasciitilde10.5 降至
6.0--8.0），参与酸碱催化或共价催化。例如：acetoacetate decarboxylase 中
Lys-115 的 pKa ≈ 6.0 {[}Bartlett et al., \emph{J. Mol. Biol.} 324, 105
(2002){]}；ribonuclease A 中 Lys-41 的 pKa ≈ 8.8 {[}Raines, \emph{Chem.
Rev.} 98, 1045 (1998){]}。高度保守，对突变极度敏感（ΔΔG \textgreater{} 3
kcal/mol 常见）。 - \textbf{Surface loop Lys-289}：溶剂暴露，pKa
接近溶液值（10.4--10.6），功能上多参与非特异性静电相互作用或蛋白质-蛋白质识别界面，保守性低，对突变容忍度高（ΔΔG
通常 \textless{} 1 kcal/mol）。

\textbf{Physical Judgment}：

\[I(Y; P \mid S) > 0 \quad \text{——在给定"氨基酸类型 = Lys"的Conditions下，残基Location携带大量功能信息}\]

\begin{longtable}[]{@{}
  >{\raggedright\arraybackslash}p{(\linewidth - 6\tabcolsep) * \real{0.1311}}
  >{\raggedright\arraybackslash}p{(\linewidth - 6\tabcolsep) * \real{0.3443}}
  >{\raggedright\arraybackslash}p{(\linewidth - 6\tabcolsep) * \real{0.3607}}
  >{\raggedright\arraybackslash}p{(\linewidth - 6\tabcolsep) * \real{0.1639}}@{}}
\toprule\noalign{}
\begin{minipage}[b]{\linewidth}\raggedright
Physical Quantity
\end{minipage} & \begin{minipage}[b]{\linewidth}\raggedright
Lys-37（Active Site）
\end{minipage} & \begin{minipage}[b]{\linewidth}\raggedright
Lys-289（Surface loop）
\end{minipage} & \begin{minipage}[b]{\linewidth}\raggedright
Difference来源
\end{minipage} \\
\midrule\noalign{}
\endhead
\bottomrule\noalign{}
\endlastfoot
侧链 pKa & 6.0--8.0 & 10.4--10.6 & 局部静电环境（Buried vs 溶剂暴露） \\
SASA (溶剂可及Surface积) & \textless{} 10 Ų & \textgreater{} 100 Ų & 3D
折叠Location \\
Evolution保守性（BLOSUM62 Location熵） & \textless{} 0.2 bits & \textgreater{}
0.8 bits & 功能约束强度 \\
突变 ΔΔG & \textgreater{} 3 kcal/mol & \textless{} 1 kcal/mol &
折叠Stability贡献 \\
\end{longtable}

\textbf{BUT------关键限制}：

Situs 的正弦编码编码的Yes\textbf{序列Location}（1D 残基编号 37 vs
289），而非\textbf{3D 空间Location}（Active Site口袋 vs Surface）。这Yes Situs
在蛋白质应用中的根本缺陷：

\begin{itemize}
\tightlist
\item
  序列Location → 3D Location映射\textbf{不YesFunction关系}。同源蛋白中，序列Location 37
  在 homolog A 中可能YesActive Site，在 homolog B 中可能Yes
  loop（插入/删除导致序列位移）。
\item
  对于没有 3D Structure的情况（大多数蛋白质预测任务），Situs
  只能使用序列Location。序列LocationYes 3D 功能Location的\textbf{弱代理变量}。
\item
  只有当我们有 3D Structure且将 Situs 应用于 3D
  坐标编码时，上述功能Difference才能被有效捕获。
\end{itemize}

\textbf{诚实Conclusion}：1D 序列正弦编码对区分 Lys-37 Active Site vs Lys-289 Surface
loop \textbf{有信息但信号弱}。序列Location与 3D 功能Location的相关性受限于 fold
保守性。\(I(Y; P_{\text{1D}} \mid S) > 0\) 但远小于
\(I(Y; P_{\text{3D}} \mid S)\)。

\begin{center}\rule{0.5\linewidth}{0.5pt}\end{center}

\subsubsection{1.1.2 α-螺旋的 3.6
残基/圈周期性}\label{ux3b1-ux87baux65cbux7684-3.6-ux6b8bux57faux5708ux5468ux671fux6027}

\textbf{Physical Facts}：

Pauling 经典 α-螺旋 {[}Pauling et al., \emph{PNAS} 37, 205 (1951){]}： -
每圈 3.6 个残基（精确值 = 3.6，或 18 残基/5 圈) - 螺距 = 5.4
Å（每残基上升 1.5 Å) - 氢键Mode：残基 \(i\) 的 C=O ··· H--N 残基
\(i+4\) - 这意味着残基 \(i\) 和 \(i+4\) 在空间上相邻（\textasciitilde3.1
Å），而 \(i\) 和 \(i+3\)（或 \(i+5\)）在螺旋的相对面

\textbf{正弦编码能No捕获？}

对于正弦编码 \(\sin(2\pi p / \lambda_j)\)： - 若波长
\(\lambda_j = 3.6\)：\(\sin(2\pi p/3.6)\) 确实具有 3.6 残基的精确周期 -
残基 \(i\) 和 \(i+3.6\) 获得相同编码值 → 残基 \(i\) 和 \(i+4\)（距离
4，周期差 0.4/3.6 ≈ 11\%）编码相近但不相同

\textbf{但Yes------标准 Transformer 编码做不到}：

标准正弦编码 {[}Vaswani et al., \emph{NeurIPS} 2017{]} 的波长集合：
\[\lambda_j = 2\pi \cdot 10000^{2j/d}\]

当 \(d = 512\) 时，波长范围约为 \(2\pi\) 到 \(2\pi \cdot 10000\)（≈ 6.28
到 ≈ 62832）。对于典型的蛋白质长度 \(L \sim 300\)： -
最短波长（\(j = d/2-1 = 255\)):\(\lambda_{255} = 2\pi \cdot 10000^{510/512} \approx 2\pi \cdot 10000^{0.996} \approx 6.28 \times 9550 \approx 60,000\)
------远大于 3.6！ - 没有任何一个 \(\lambda_j\) 接近 3.6

\textbf{PhysicsOptimal谱（定理 1.2.1）能No解决？}

根据定理 1.2.1，Optimal频率谱需要知道Physics相关长度 \(\xi\)。对于
α-螺旋，\(\xi \approx 3.6\)（螺旋周期）。若我们以 \(\xi = 3.6\)
为Objective设计编码： -
\(\lambda_{\min} = 2\pi\xi \cdot \cot(\pi(2(d/2-1)+1)/2d) \approx 2\pi \cdot 3.6 \cdot \cot(\pi/2 - \pi/2d) \approx 22.6 \cdot (2d/\pi) \approx 14.4 d\)
- 对于 \(d = 64\)，\(\lambda_{\min} \approx 920\) ------仍远大于 3.6！

\textbf{根本Problem}：正弦编码的''分辨率''受限于 \(d\)。要在 \(L = 300\)
的序列上分辨出 3.6 残基的周期，需要
\(\lambda_{\min} \lesssim 3.6\)，即需要
\(d \gtrsim 300/3.6 \times 2 = 167\) 维（每个周期至少 2
个采样点，Nyquist）。对于典型的 \(d = 64\) 或 \(d = 128\)，3.6
残基的周期性\textbf{Theory上None法被捕获}。

\textbf{诚实Conclusion}：对于实际使用的编码Dimension（\(d \leq 128\)），正弦编码\textbf{None法}可靠捕获
α-螺旋的 3.6
残基/圈周期性。这不Yes编码形式的缺陷，而YesDimension-分辨率的基本权衡（Corollaries
1.2.1）。如果需要在序列Location编码中捕获螺旋周期性，应该显式添加一个
\(\lambda \approx 3.6\) 的频率分量，或者使用可Learning的周期性Characteristics。

\begin{center}\rule{0.5\linewidth}{0.5pt}\end{center}

\subsubsection{1.1.3 固有None序区域（IDR）中 Situs
的退化}\label{ux56faux6709ux65e0ux5e8fux533aux57dfidrux4e2d-situs-ux7684ux9000ux5316}

\textbf{Physical Facts}：

固有None序区域（IDR）占真核蛋白质组的 30--40\% {[}Dunker et al.,
\emph{FEBS J.} 272, 5129 (2005){]}。关键Characteristics： -
在生理Conditions下\textbf{不折叠为单一稳定 3D Structure} -
功能来自\textbf{氨基酸组成}（如电荷Mode、脯氨酸含量）而非精确折叠 -
序列Location与功能的关系\textbf{极弱}：IDR 的功能（如 linker、entropic
spring、PTM 位点集群）主要由组成决定，Location排列有高度容忍度 - 同一 IDR
插入/删除 5--20 个残基通常不影响功能 {[}Zibaee et al., \emph{Protein
Sci.} 16, 906 (2007){]}

\textbf{Situs 在 IDR 上的lines为}：

\[I(Y; P_{\text{seq}} \mid S = \text{"IDR residue"}) \approx 0\]

由定理 2.2.1 和 2.3.1：
\[\delta_s^{\text{PE}} \leq \sqrt{\frac{1}{2} I(Y; P \mid S)} + \delta_s^{\text{variance}} \approx 0 + O(1/\sqrt{M})\]

即：\textbf{Situs 退化为Dimension为 \(d\) 的加法Noise}。

具体地，在 IDR 区域的 \(L\) 个连续Location上：
\[\text{PE}(p) = [\sin(2\pi p/\lambda_1), \cos(2\pi p/\lambda_1), \ldots]\]

这些编码向量随 \(p\) 平滑变化（Lipschitz
连续性），但它们携带的变异性与标签 \(Y\) \textbf{None关}。在Training中： -
CleanSample：\(\text{PE}(p)\) 添加了与标签None关的波动 →
可能略微降低Expert consensus →
\(p_{\text{clean}}^{\text{PPE}} \geq p_{\text{clean}}\) -
NoiseSample：\(\text{PE}(p)\) 同样不提供区分Noise的能力 →
\(p_{\text{noisy}}^{\text{PPE}} \approx p_{\text{noisy}}\) -
Result：\(\delta_s^{\text{PE}} \leq 0\)（负或零）

\textbf{诚实Conclusion}：IDR Yes Situs
的\textbf{最坏场景}------Location编码退化为None害但None用的Dimension膨胀。它不会主动破坏Model（Lipschitz
连续性Guarantee \(h_i\) 仍然光滑），但消耗了 \(d\)
维的Computation和存储，且可能通过梯度Noise略微降低Convergence速度。如果Data集中 IDR
占比高（\textgreater30\% 如真核蛋白），Situs 的净收益\textbf{可能为负}。

\begin{center}\rule{0.5\linewidth}{0.5pt}\end{center}

\subsection{1.2 Materials 3D
StructureLocation编码}\label{ux6750ux6599-3d-ux7ed3ux6784ux4f4dux7f6eux7f16ux7801}

\subsubsection{1.2.1 V\_N 空位形成能：晶界 vs 体相 vs
Surface}\label{v_n-ux7a7aux4f4dux5f62ux6210ux80fdux6676ux754c-vs-ux4f53ux76f8-vs-ux8868ux9762}

\textbf{Physical Facts（DFT ComputationData）}：

AlN 中氮空位 \(V_N\) 的形成能强烈Dependencies于空间Location {[}Freysoldt et al.,
\emph{Rev.~Mod. Phys.} 86, 253 (2014); Stampfl \& Van de Walle,
\emph{Phys. Rev.~B} 65, 155212 (2002){]}：

\begin{longtable}[]{@{}
  >{\raggedright\arraybackslash}p{(\linewidth - 4\tabcolsep) * \real{0.1714}}
  >{\raggedright\arraybackslash}p{(\linewidth - 4\tabcolsep) * \real{0.5429}}
  >{\raggedright\arraybackslash}p{(\linewidth - 4\tabcolsep) * \real{0.2857}}@{}}
\toprule\noalign{}
\begin{minipage}[b]{\linewidth}\raggedright
Location
\end{minipage} & \begin{minipage}[b]{\linewidth}\raggedright
形成能 \(E_f\) (eV)
\end{minipage} & \begin{minipage}[b]{\linewidth}\raggedright
PhysicsReason
\end{minipage} \\
\midrule\noalign{}
\endhead
\bottomrule\noalign{}
\endlastfoot
\textbf{体相}（bulk, 配位数 4） & 3.2--4.0 & Complete四面体配位，断开 4 个
Al--N 键 \\
\textbf{Surface}（(0001) Surface, 配位数 3） & 1.8--2.5 &
Surface弛豫释放应变能，配位数降低 \\
\textbf{晶界}（Σ7 对称倾斜晶界） & 1.0--1.8 & 晶界本征应力场 + 多余体积
+ 悬挂键预存在 \\
\end{longtable}

Difference幅度：\(\Delta E_f \approx 1.5\)--\(2.5\) eV，远超 \(k_B T\)（室温下
≈ 0.025 eV），在\textbf{任何}温度下都具有决定性影响。

\textbf{Physics因果链}（不可简化):

\[\text{3D Location} \xrightarrow{\text{决定}} \text{局部配位环境} \xrightarrow{\text{决定}} \text{缺陷形成能}\]

这is a清晰的因果箭头------LocationYes自变量，形成能Yes因变量。Location编码捕获的Yes这个因果链的\textbf{Reason侧}。

\[I(Y = E_f; P_{\text{3D}} \mid S = V_N) > 0 \quad \text{——Information Theory意义上严格正且量大}\]

根据 Fano Lower Bound（Theorem 2.4.1），对于多类分类（不同缺陷电荷态):
\[\delta_s^{\text{PE}} \geq \frac{2 \cdot I(Y; P \mid S) - \log 2}{\log |\mathcal{Y}|}\]

取 \(I(Y; P \mid S) \approx H(Y|S) \approx \log 3 \approx 1.6\) bits（3
种形成能区间：高/中/低），\(\log |\mathcal{Y}| \approx \log 3\)：
\[\delta_s^{\text{PE}} \gtrsim 0.5 \text{–}0.8\]

这is a\textbf{可观的Detection边际增益}，对应 Theorem 1 的指数界显著收紧。

\textbf{诚实Conclusion}：对于Materials缺陷Detection，3D Situs
Yes\textbf{Physics上最有价值的应用场景}。Location直接因果决定标签。这Yes Situs
的''home turf''。

\begin{center}\rule{0.5\linewidth}{0.5pt}\end{center}

\subsubsection{1.2.2 3D
旋转编码的Physics对称性检查}\label{d-ux65cbux8f6cux7f16ux7801ux7684ux7269ux7406ux5bf9ux79f0ux6027ux68c0ux67e5}

\textbf{Situs 旋转编码的定义}（定义 1.3.1):
\[\text{PE}_{\text{rot}}(\mathbf{p}) = \mathbf{R}_x(\alpha x) \cdot \mathbf{R}_y(\beta y) \cdot \mathbf{R}_z(\gamma z) \cdot \mathbf{e}_0\]

\textbf{检查 1：平移不变性} ✓

内积：
\[\langle \text{PE}(\mathbf{p}), \text{PE}(\mathbf{q}) \rangle = \langle \mathbf{e}_0, \mathbf{R}(\mathbf{q} - \mathbf{p}) \mathbf{e}_0 \rangle = \cos(\alpha\Delta x) + \cos(\beta\Delta y) + \cos(\gamma\Delta z)\]

仅Dependencies于
\(\Delta\mathbf{p} = \mathbf{q} - \mathbf{p}\)。\textbf{平移不变性严格成立}。这满足晶体学的基本要求------完美晶体中，两个Location等价当且仅当它们的相对位移Yes晶格向量。

\textbf{检查 2：旋转等变性} ✗ \textbf{不成立}

考虑全局旋转 \(R_{\text{phys}} \in SO(3)\)
作用于坐标：\(\mathbf{p} \to R_{\text{phys}} \mathbf{p}\)。

对于 \$R\_\{\text{phys}\} = \$ 绕 z 轴旋转 90°（将 \((x, y, z)\) 变为
\((-y, x, z)\)):
\[\text{PE}_{\text{rot}}(R_{\text{phys}}\mathbf{p}) = \mathbf{R}_x(-\alpha y) \mathbf{R}_y(\beta x) \mathbf{R}_z(\gamma z) \mathbf{e}_0\]

而
\(\text{PE}_{\text{rot}}(\mathbf{p}) = \mathbf{R}_x(\alpha x) \mathbf{R}_y(\beta y) \mathbf{R}_z(\gamma z) \mathbf{e}_0\)。

这两个编码之间\textbf{不存在}一个与 \(\mathbf{p}\) None关的变换 \(T\) 使得
\(T \cdot \text{PE}(\mathbf{p}) = \text{PE}(R_{\text{phys}}\mathbf{p})\)。Reason：编码为三个坐标轴分配了\textbf{不相交的Dimension对}------\((0,1)\)
对应 \(x\)，\((2,3)\) 对应 \(y\)，\((4,5)\) 对应 \(z\)。Physics旋转混合了
\(x\) 和 \(y\) 分量，但编码空间中没有对应的''混合''Mechanism。

\textbf{破坏的后果}：

在以下场景中旋转等变性的缺失Yes致命的： 1.
\textbf{多晶/晶界System}：不同晶粒的取向不同，相同类型缺陷的编码因取向不同而不同
→ \(h_i\) 引入与标签None关的变异性 2.
\textbf{分子动力学轨迹}：分子整体旋转导致所有原子编码变化 →
Model可能学到''旋转Status''而非''PhysicsStatus'' 3. \textbf{Surface slab
Computation}：同一Surface不同切面（如 AlN (0001) vs (11̄00)）具有不同的坐标轴取向

\textbf{可以修复吗？}

Theory上可以。替代方案： -
\textbf{球谐Function编码}：\(\text{PE}(\mathbf{p}) = [Y_{\ell m}(\theta, \phi)]\)，具有严格的
\(SO(3)\) 旋转等变性 - \textbf{E(3) 等变网络}：使用 Tensor Field
Networks {[}Thomas et al., \emph{NeurIPS} 2018{]} 或 SE(3)-Transformers
{[}Fuchs et al., \emph{NeurIPS} 2020{]}，天然具备平移+旋转等变性 -
\textbf{Cost}：球谐Function编码的Dimension随角动量 \(L\) 增长为
\((L+1)^2\)，\(L=3\) 即需 16 维（vs 6 维的旋转编码），且 Lipschitz
lines为更复杂

\textbf{诚实Conclusion}：Situs 的 3D
旋转编码\textbf{仅具有平移不变性，不具旋转等变性}。对于大多数Materials应用（需要处理任意取向的晶体），这Yes\textbf{严重缺陷}。但在受控场景（单晶、固定取向的
slab、取向已对齐的Data集）中，这不构成Problem。

\begin{center}\rule{0.5\linewidth}{0.5pt}\end{center}

\subsection{1.3 Total Atom Count N: 32 vs 256 Supercell DFT
Accuracy}\label{ux539fux5b50ux603bux6570-n32-vs-256-ux8d85ux80deux7684-dft-ux7cbeux5ea6}

\subsubsection{1.3.1
System性Difference的Physics起源}\label{ux7cfbux7edfux6027ux5deeux5f02ux7684ux7269ux7406ux8d77ux6e90}

\textbf{Makov-Payne 有限尺寸标度} {[}Makov \& Payne, \emph{Phys. Rev.~B}
51, 4014 (1995){]}：

带电缺陷的形成能随超胞尺寸 \(L\) 的标度lines为：
\[E_f(L) = E_f(\infty) + \frac{\alpha q^2}{2\varepsilon L} + \frac{2\pi q Q}{3\varepsilon L^3} + O(L^{-5})\]

where \(q\) Yes缺陷电荷，\(\varepsilon\) Yes介电常数，\(Q\) Yes弹性四极矩。

\textbf{具体数值估计（\(V_N^{+1}\) 在 AlN
中，\(\varepsilon \approx 9\)）}：

\begin{longtable}[]{@{}
  >{\raggedright\arraybackslash}p{(\linewidth - 10\tabcolsep) * \real{0.1538}}
  >{\raggedright\arraybackslash}p{(\linewidth - 10\tabcolsep) * \real{0.1231}}
  >{\raggedright\arraybackslash}p{(\linewidth - 10\tabcolsep) * \real{0.2308}}
  >{\raggedright\arraybackslash}p{(\linewidth - 10\tabcolsep) * \real{0.1692}}
  >{\raggedright\arraybackslash}p{(\linewidth - 10\tabcolsep) * \real{0.2000}}
  >{\raggedright\arraybackslash}p{(\linewidth - 10\tabcolsep) * \real{0.1231}}@{}}
\toprule\noalign{}
\begin{minipage}[b]{\linewidth}\raggedright
超胞尺寸
\end{minipage} & \begin{minipage}[b]{\linewidth}\raggedright
原子数
\end{minipage} & \begin{minipage}[b]{\linewidth}\raggedright
周期性图像间距
\end{minipage} & \begin{minipage}[b]{\linewidth}\raggedright
\(1/L\) 修正
\end{minipage} & \begin{minipage}[b]{\linewidth}\raggedright
\(1/L^3\) 修正
\end{minipage} & \begin{minipage}[b]{\linewidth}\raggedright
Total Correction
\end{minipage} \\
\midrule\noalign{}
\endhead
\bottomrule\noalign{}
\endlastfoot
2×2×2 & 32 & \textasciitilde6.2 Å & \textasciitilde0.35 eV &
\textasciitilde0.12 eV & \textbf{\textasciitilde0.5 eV} \\
2×2×3 & 48 & \textasciitilde6.2--9.3 Å & \textasciitilde0.25 eV &
\textasciitilde0.05 eV & \textbf{\textasciitilde0.3 eV} \\
3×3×3 & 108 & \textasciitilde9.3 Å & \textasciitilde0.15 eV &
\textasciitilde0.02 eV & \textbf{\textasciitilde0.17 eV} \\
4×4×4 & 256 & \textasciitilde12.4 Å & \textasciitilde0.08 eV &
\textasciitilde0.005 eV & \textbf{\textasciitilde0.09 eV} \\
\end{longtable}

\textbf{PhysicsOrigin}（三个独立的有限尺寸效应，均不可忽略):

\begin{enumerate}
\def\labelenumi{\arabic{enumi}.}
\tightlist
\item
  \textbf{静电像电荷相互作用}（\(1/L\) 主导项):

  \begin{itemize}
  \tightlist
  \item
    带电缺陷与其周期性镜像之间的 Coulomb 相互作用
  \item
    32 原子超胞中相邻镜像仅隔 6.2 Å → 相互作用能 \textasciitilde0.3--0.5
    eV
  \item
    Physics上这Yes人为引入的------真实孤立缺陷没有周期性镜像
  \end{itemize}
\item
  \textbf{弹性像相互作用}（\(1/L^3\)):

  \begin{itemize}
  \tightlist
  \item
    缺陷引起的局部应变场与其周期性镜像的相互作用
  \item
    对 \(V_N\)（移除原子产生拉伸应变）此效应不可忽略
  \end{itemize}
\item
  \textbf{缺陷波Function重叠}（指数衰减 \(\propto e^{-L/\xi}\)):

  \begin{itemize}
  \tightlist
  \item
    对深能级缺陷（如 \(V_N\) 在 AlN 中产生 \textasciitilde1.5 eV
    深的施主能级），波Function局域化半径 \(\xi \approx 2\)--\(3\) Å
  \item
    32 原子超胞中相邻波Function尾重叠不可忽略 →
    缺陷能级分裂成能带（人造色散）
  \end{itemize}
\end{enumerate}

\textbf{这与 Situs 的 \(N\) 编码有什么关系？}

Situs 的Part三域（定义 1.1.3）将总原子数 \(N\)
作为标量Location编码入。ProblemYes：

\begin{itemize}
\tightlist
\item
  \(N=32\) 和 \(N=256\)
  的Data\textbf{不Yes同一PhysicsSystem的不同''Location''}------它们Yes\textbf{不同Accuracy级别}的近似
\item
  \(N\) 编码将''ComputationAccuracy''与''Physics标签''混淆。\(N=32\) 的 \(E_f\)
  偏差YesComputation artifact，不YesPhysics信号
\item
  如果Training集混合了 32 和 256 原子超胞的Data，\(N\)
  编码可能学到''\(N=32 \to\) 高 \(E_f\)``的虚假相关性
\end{itemize}

\textbf{Physics上Correct的处理}： - 对 \(E_f\) 应用 Freysoldt-Neugebauer-Van
de Walle (FNV) 有限尺寸修正 {[}Freysoldt et al., \emph{Phys. Rev.~Lett.}
102, 016402 (2009){]}，消除 \(N\) 的System性偏差 -
\textbf{之后}再考虑YesNo使用 \(N\) 编码------此时 \(N\)
才反映真实的System尺寸效应（如量子限域），而非Computation artifact

\textbf{诚实Conclusion}：如果TrainingData未经过有限尺寸修正，\(N\)
编码学到的Yes\textbf{Computation artifact 而非Physics}。在修正之后，\(N\)
携带的真实Physics信息（量子限域、振动Mode数）对于缺陷Physics来说较小（\(\Delta E_f\)
的 \(N\)-dependence \textless{} 0.05 eV for
\(N > 100\)）。对于大多数Materials任务，\(N\) 编码的 \(\delta_s^{\text{PE}}\)
\textbf{接近零}------这is aPhysics上动机弱、实践上危险（容易学到
artifact）的编码Dimension。

\begin{center}\rule{0.5\linewidth}{0.5pt}\end{center}

\section{Part 2 部分：Situs
在什么情况下有害？}\label{ux7b2c-2-ux90e8ux5206situs-ux5728ux4ec0ux4e48ux60c5ux51b5ux4e0bux6709ux5bb3}

诚实暴击：Situs
不is a''None害但可能None用''的组件。在特定Conditions下它Yes\textbf{主动有害的}。

\begin{center}\rule{0.5\linewidth}{0.5pt}\end{center}

\subsection{2.1 虚假相关性（Spurious
Correlation）}\label{ux865aux5047ux76f8ux5173ux6027spurious-correlation}

\textbf{Mechanism}：

当PhysicsLocation \(P\) 与标签 \(Y\)
在TrainingData中具有统计学相关，但该相关\textbf{不反映因果关系}时：

\[\text{Corr}_{\text{train}}(P, Y) \neq 0, \quad \text{但 } P \perp\!\!\!\perp Y \text{（因果意义上独立）}\]

由于 Situs 将 \(P\) 编码为 \(\text{PE}(P)\) 并注入表示，Model（专家
\(E_m\)）会Learning利用这个虚假相关。

\textbf{具体实例}：

\textbf{蛋白质Data集偏差}： - PDB
中解析的蛋白质Structure存在''易于结晶''偏差------倾向于更稳定、更紧凑的折叠 -
如果Training集中所有酶Active Site恰好在结晶Structure中位于蛋白质核心（低 B-factor
区域 → 低Location索引附近被解析得更清楚) - Situs
学到：残基Location靠前（\(p < 100\)）→ Active Site - Test集中出现 C
端Active Site蛋白 → System性误判

\textbf{DFT ComputationParameter偏差}： - Materials Project
中的Computation使用不同的ComputationParameter（赝势、k
点密度），且这些Parameter与Materials类型有System性关联 - 如果 Situs 的原子数编码
\(N\) 学到：\(N=32\)（典型金属Computation）→
低形成能，\(N=256\)（典型半导体缺陷Computation）→ 高形成能 -
这种''知识''YesComputation惯例的 artifact，不YesPhysics

\textbf{虚假相关性的Detection}：\(\delta_s^{\text{PE}}\) 在Training集上
\textgreater{} 0，但在 i.i.d. 验证集上 → 0 或 \textless{} 0。Theorem
2'（修正 Theorem 2）的Upper Bound在Distribution外Data上\textbf{反而更松}------因为
\(\varepsilon_{\text{PE}}\) 在TrainingDistribution上被低估。

\textbf{诚实Conclusion}：虚假相关性Yes Situs
的\textbf{最大实际风险}。Physics学家知道''position matters''但TrainingData中的
position-label
相关有多少来自因果、有多少来自Data收集偏差，必须逐案Analysis。\textbf{没有捷径}。

\begin{center}\rule{0.5\linewidth}{0.5pt}\end{center}

\subsection{2.2
多重共线性（Multicollinearity）}\label{ux591aux91cdux5171ux7ebfux6027multicollinearity}

\textbf{Mechanism}：

当Status表示 \(\phi(s_i)\) \textbf{已经隐式编码了Location信息}时，加法注入
\(\text{PE}(p_i)\) 产生冗余Dimension：

\[h_i = \underbrace{\phi(s_i)}_{\text{已含隐式Location信息}} + \underbrace{\text{PE}(p_i)}_{\text{显式Location信息}}\]

这导致编码空间中的有效Dimension膨胀，伴随： 1.
\textbf{专家Model的梯度方差增大}：每个专家需要From冗余Dimension中提取信号 2.
\textbf{有效专家多样性 \(\rho_{\text{eff}}\)
下降}：当所有专家在同一冗余空间中操作，它们的输出相关性增大

\textbf{具体实例}：

\textbf{图神经网络 (GNN) 处理晶体}： - GNN
已通过消息传递聚合了邻域信息，包含隐式的Location/距离编码 -
如果再对每个节点添加 Situs 的 3D
旋转编码，相当于添加了\textbf{已经被图卷积层捕获的信息} -
Result：\(I(Y; \text{PE}(P) \mid \phi(S)) \approx 0\)（给定 GNN
Characteristics后，LocationNone额外信息) - 但 PE 的额外Dimension增加了优化难度，等效于
\(\delta_s^{\text{variance}}\) 增大

\textbf{蛋白质语言Model (pLM) 处理序列}： - ESM-2 {[}Lin et al.,
\emph{Science} 379, 1123 (2023){]} 等 pLM 的 attention
Mechanism天然包含Location信息（attention bias) - 对 ESM embedding 再加 Situs
正弦编码 = 添加已被 attention 处理的冗余信号 -
这种''双重编码''在统计上引入共线性，但不增加有效信息

\textbf{定量影响}：

设原始Characteristics空间的有效秩为
\(r_{\text{eff}} = \text{rank}(\Sigma_\phi)\)。添加 PE 后：
\[\text{rank}(\Sigma_h) \leq r_{\text{eff}} + \min(d_{\text{pe}}, N)\]

但如果 \(\text{span}(\text{PE}(P)) \subseteq \text{span}(\phi(S))\)（PE
在 \(\phi\) 的列空间中），则 \(\text{rank}(\Sigma_h) = r_{\text{eff}}\)
------Dimension增加了但有效秩不变。这Yes\textbf{纯浪费}。

\textbf{诚实Conclusion}：在已有强大Structure编码的Model（GNN、Transformer、pLM）之上叠加
Situs，大概率Yes\textbf{零增益 + 微损害}。Situs
最有用时Yes作为\textbf{仅有的}Location信息源（当 \(\phi(s_i)\)
完全没有空间信息时）。

\begin{center}\rule{0.5\linewidth}{0.5pt}\end{center}

\subsection{2.3 Physics对称性破坏（Physical Symmetry
Breaking）}\label{ux7269ux7406ux5bf9ux79f0ux6027ux7834ux574fphysical-symmetry-breaking}

\textbf{Mechanism}：

Situs 的加法注入 \(h_i = \phi(s_i) + \text{PE}(p_i)\)
\textbf{刻意破坏}了PhysicsSystem的某些对称性。当这些对称性对预测任务重要时，这Yes有害的。

\textbf{三种破坏的对称性}：

\subsubsection{2.3.1
平移对称性破坏（对体相性质预测Yes坏的）}\label{ux5e73ux79fbux5bf9ux79f0ux6027ux7834ux574fux5bf9ux4f53ux76f8ux6027ux8d28ux9884ux6d4bux662fux574fux7684}

完美晶体中，所有原胞等价------体相性质（带隙、弹性常数、声子频率）应具有\textbf{完全平移不变性}。

使用 Situs 后： - 两个等价原胞中的原子获得不同的 \(h_i\) -
Model必须额外Learning''忽略Location''------增加了不必要的Learning负担 -
更糟：如果Training集有限，Model可能学到''带隙Dependencies于原子索引''的虚假规律

\textbf{实例}：预测 AlN 体相带隙（\textasciitilde6.2 eV，{[}Yan et al.,
\emph{Phys. Rev.~B} 93, 014110 (2016){]}) - \(N=32\) 超胞中每个 N
原子的 Situs 编码不同 -
如果Model用这些编码预测同一个带隙标签，它必须在内部Learning''这些编码应该被忽略''
- 对小Data集，这可能导致过拟合到特定的编码值

\textbf{判断}：如果任务标签具有平移不变性（体相性质），Situs
\textbf{不应使用}。如果任务标签不具有平移不变性（缺陷形成能、Surface能、局域态密度），Situs
\textbf{应当使用}。

\subsubsection{2.3.2 旋转对称性破坏（3D 旋转编码的硬伤------已在 1.2.2
Analysis）}\label{ux65cbux8f6cux5bf9ux79f0ux6027ux7834ux574f3d-ux65cbux8f6cux7f16ux7801ux7684ux786cux4f24ux5df2ux5728-1.2.2-ux5206ux6790}

\subsubsection{2.3.3
排列对称性破坏（对集合型预测Yes坏的）}\label{ux6392ux5217ux5bf9ux79f0ux6027ux7834ux574fux5bf9ux96c6ux5408ux578bux9884ux6d4bux662fux574fux7684}

原子System本质上Yes\textbf{粒子的集合}------原子的编号/顺序Yes人为的，不应影响Physics预测。

Situs Dependencies''Part \(i\) 个原子的Location \(p_i\)``，where \(i\)
Yes原子的索引。如果Data加载时的原子排序发生变化，\(h_i\)
会不同------即使PhysicsSystem完全相同。

\textbf{这在以下场景中造成Problem}： - DFT
Computation中原子列表的排序NonePhysics意义（VASP/PWscf
输出中的原子顺序Yes输入顺序的Function) -
分子动力学轨迹中，原子可能跨过周期性边界，改变其在数组中的Location -
不同Data源的原子排序约定不同

\textbf{缓解措施}：使用\textbf{等变}的Location编码（如基于距离矩阵或原子环境的编码），或要求Data预处理中Guarantee一致的原子排序（如按坐标排序------但这又引入了排序的人为性）。

\textbf{诚实Conclusion}：Situs
破坏排列对称性is a\textbf{实际工程Problem}而非Theory缺陷------可以通过一致的预处理缓解，但None法根除。如果原子排序在不同Sample间不一致，Situs
引入的Noise可能淹没信号。

\begin{center}\rule{0.5\linewidth}{0.5pt}\end{center}

\section{Part 3
部分：适用场景Analysis表}\label{ux7b2c-3-ux90e8ux5206ux9002ux7528ux573aux666fux5206ux6790ux8868}

对 6 个具体场景，Evaluation \(\delta_s^{\text{PE}}\)
的Symbol（正/零/负）、量级（大/中/小/零）和物Theory证。

\begin{center}\rule{0.5\linewidth}{0.5pt}\end{center}

\begin{longtable}[]{@{}
  >{\raggedright\arraybackslash}p{(\linewidth - 8\tabcolsep) * \real{0.0577}}
  >{\raggedright\arraybackslash}p{(\linewidth - 8\tabcolsep) * \real{0.1154}}
  >{\raggedright\arraybackslash}p{(\linewidth - 8\tabcolsep) * \real{0.5192}}
  >{\raggedright\arraybackslash}p{(\linewidth - 8\tabcolsep) * \real{0.1923}}
  >{\raggedright\arraybackslash}p{(\linewidth - 8\tabcolsep) * \real{0.1154}}@{}}
\toprule\noalign{}
\begin{minipage}[b]{\linewidth}\raggedright
\#
\end{minipage} & \begin{minipage}[b]{\linewidth}\raggedright
场景
\end{minipage} & \begin{minipage}[b]{\linewidth}\raggedright
\(\delta_s^{\text{PE}}\) Symbol
\end{minipage} & \begin{minipage}[b]{\linewidth}\raggedright
量级估计
\end{minipage} & \begin{minipage}[b]{\linewidth}\raggedright
理由
\end{minipage} \\
\midrule\noalign{}
\endhead
\bottomrule\noalign{}
\endlastfoot
\textbf{1} & \textbf{AlN 缺陷Detection：\(V_N\) 空位，534 帧} & \textbf{正} &
\textbf{大} & 3D Location因果决定缺陷形成能（§1.2.1）。534
帧在Materials缺陷Data中属于中等规模，足以LearningLocation-能量关系。Location信息不冗余（\(\phi(s_i)\)
为局部ChemistryCharacteristics，不含全局空间上下文）。3D
旋转编码的平移不变性在此场景中恰好合适（缺陷中心对齐后相对Location有意义）。\(I(Y; P \mid S)\)
估计：1.5--2.5 bits（通过形成能三分位
binning），\(\delta_s^{\text{PE}} \gtrsim 0.4\)。\textbf{但Yes}：如果 534
帧包含多个晶粒取向，旋转等变性缺失会引入Noise（见 §1.2.2
的警告）。建议先检查取向Distribution，必要时做取向对齐预处理。 \\
\textbf{2} & \textbf{蛋白质功能预测：酶Active Site分类} & \textbf{正}（弱）
& \textbf{中小} & 序列Location与功能位点有统计相关性（N 端/C
端偏好、Structure域Location），但效应量弱。\(I(Y; P_{\text{1D}} \mid S) > 0\)
成立但数值有限（\textless{} 0.3
bits，因为大部分功能信息已在残基类型和局部序列上下文中）。\textbf{关键限制}：Situs
使用 1D 序列Location而非 3D 空间Location------而Active Site功能由 3D
Location决定。序列→3D 映射因 fold
Difference而非单值。\textbf{Conclusion}：有微弱正面贡献，但远不如直接使用 3D
Structure信息（如有 AF2 预测Structure）。若已有 pLM
embedding（含Location信息），收益趋近于零（共线性Problem，§2.2）。 \\
\textbf{3} & \textbf{Drug-target 对接评分：DrugBank 6784 条} &
\textbf{正}（有Conditions） & \textbf{中小} & conditions:使用 3D 对接姿态坐标（而非
1D 序列Location）。结合亲和力（\(K_d\), \(K_i\), \(\Delta G\)）Yes 3D
姿态的Function------口袋Location、氢键几何、疏水接触面积都由空间Location决定。\(I(Y; P_{\text{3D}} \mid S)\)
\textgreater{} 0。\textbf{但Yes}：6784 条Data对于Learning复杂的 3D
Location-亲和力映射偏少（特别Yes考虑到对接姿态的巨大构象空间）。Situs 的
Lipschitz 光滑性（定理
2.5.1）在此Yes有益的------它正则化了邻近空间Location的预测。有过度Dependencies结合位点Location的风险（分子骨架Structure也有贡献）。\textbf{建议}：仅在对接姿态质量高（RMSD
\textless{} 2.0 Å to experimental）的子集上使用 3D Situs。 \\
\textbf{4} & \textbf{纯Chemistry组成分类：None空间Structure} & \textbf{零}（或略负）
& \textbf{零} & \(\mathcal{P}\)
未定义。没有PhysicsLocation信息可编码。强lines使用（如人工赋予虚拟索引）→
编码的YesData加载顺序的 artifact。\(I(Y; P \mid S) \equiv 0\)。由定理
2.2.1，\(\delta_s^{\text{PE}} = 0\)。由定理
2.5.1，编码随虚拟Location平滑变化------但这变化不含任何标签信息。唯一的非零效应来自
\(\delta_s^{\text{variance}}\)（有限专家的估计Noise），表现为
\(O(1/\sqrt{M})\) 的\textbf{净负贡献}。\textbf{这Yes Situs
的绝对禁区}------加了一定更差。 \\
\textbf{5} & \textbf{随机排列的原子构型：None周期Structure} & \textbf{负} &
\textbf{中} & Location存在（有 3D
坐标），但Location与标签之间\textbf{没有Physics规律性}------构型Yes随机的。\(I(Y; P \mid S) = 0\)（随机排列破坏了所有Location-标签Information Theory联系）。但Training集中有限的随机涨落可能产生\textbf{偶然的}Location-标签相关，Model会学到这些虚假Mode
→ Test时泛化差。由 Theorem 2'（不完美 PPE
Lower Bound），\(\varepsilon_{\text{PE}}\)
在TrainingDistribution上被低估，导致Upper Bound过松（过度自信）。\(\delta_s^{\text{PE}}\)
在Training集上 \textgreater{} 0（虚假），在Test集上 \textless{}
0（真实）。\textbf{这Yes Situs
最危险的场景}------Surface收益Yes统计假象。使用前必须检验 \(I(Y; P \mid S)\)
的统计显著性（通过 permutation test）。 \\
\textbf{6} & \textbf{碳纳米管手性分类：(n, m) 索引} & \textbf{正} &
\textbf{大} & 手性 \((n,m)\) 完全由 3D
原子排列的全局几何决定------这YesLocation编码的理想场景。所有碳原子具有相同的局部
\(sp^2\) Chemistry环境（\(\phi(s_i) \approx\) 常数），因此
\(I(Y; S) \approx 0\)（纯Chemistry组成None法区分手性）而
\(I(Y; P) = H(Y) = \log_2 K\)（\(K\)
为手性Category数）。\(I(Y; P \mid S) \approx H(Y)\)------\textbf{几乎全部信息来自Location}。Situs
成为信息的\textbf{主要}而非辅助来源。正弦/旋转编码通过捕捉原子间的相对LocationMode，可以编码螺旋缠绕的周期性和管径信息。\textbf{但Yes}：需要注意周期性边界Conditions的编码一致性------CNT
在圆周方向有周期性，Situs 的 3D
旋转编码不会自动沿圆周闭合。对于小直径管（n+m \textless{}
15），周长效应明显，需要特殊处理。 \\
\end{longtable}

\begin{center}\rule{0.5\linewidth}{0.5pt}\end{center}

\section{Part 4
部分：每个场景一句话建议}\label{ux7b2c-4-ux90e8ux5206ux6bcfux4e2aux573aux666fux4e00ux53e5ux8bddux5efaux8bae}

\begin{longtable}[]{@{}
  >{\raggedright\arraybackslash}p{(\linewidth - 4\tabcolsep) * \real{0.0682}}
  >{\raggedright\arraybackslash}p{(\linewidth - 4\tabcolsep) * \real{0.1364}}
  >{\raggedright\arraybackslash}p{(\linewidth - 4\tabcolsep) * \real{0.7955}}@{}}
\toprule\noalign{}
\begin{minipage}[b]{\linewidth}\raggedright
\#
\end{minipage} & \begin{minipage}[b]{\linewidth}\raggedright
场景
\end{minipage} & \begin{minipage}[b]{\linewidth}\raggedright
建议（加不加 Situs + 一句话理由）
\end{minipage} \\
\midrule\noalign{}
\endhead
\bottomrule\noalign{}
\endlastfoot
\textbf{1} & AlN \(V_N\) 缺陷 534 帧 & \textbf{加}------3D
LocationYes形成能的直接PhysicsReason，534
帧足够Learning空间Mode，前提Yes取向已对齐或Model能容忍旋转非等变性。 \\
\textbf{2} & 酶Active Site分类 & \textbf{不加}（或改用 3D Situs）------1D
序列LocationYes 3D 功能Location的弱代理，尤其当已有 pLM
embedding（含注意力Location偏差）时，1D Situs 沦为共线性冗余。如有 AlphaFold
预测Structure，改用 3D Situs。 \\
\textbf{3} & DrugBank 对接评分 & \textbf{有Conditions加}------仅当使用高质量
3D 对接姿态（RMSD \textless{} 2 Å）时加 3D Situs；6784 条Data对 3D
编码Dimension偏小，需配合强正则化或降维编码。 \\
\textbf{4} & 纯组成分类NoneStructure & \textbf{绝对不加}------没有Location就没有
Situs；任何人为赋予的伪Location都Yes统计Noise，可能产生虚假相关。 \\
\textbf{5} & 随机原子构型 &
\textbf{不加}------随机Location与标签的Information Theory联系为零，Situs
在此Yes虚假相关性的滋生器；若真要使用，必须先做 permutation test 验证
\(I(Y;P \mid S)\) 的统计显著性。 \\
\textbf{6} & CNT 手性 (n,m) 分类 & \textbf{加}------手性由全局 3D
几何唯一确定，Situs
From仅有的Location信息中提取Characteristics（碳原子的局部Chemistry环境完全不携带手性信息），这Yes
Situs 的''教科书级''应用场景。 \\
\end{longtable}

\begin{center}\rule{0.5\linewidth}{0.5pt}\end{center}

\section{Summary：Situs
的Physics适用边界}\label{ux603bux7ed3situs-ux7684ux7269ux7406ux9002ux7528ux8fb9ux754c}

\subsection{它何时工作（按信心排序）}\label{ux5b83ux4f55ux65f6ux5de5ux4f5cux6309ux4fe1ux5fc3ux6392ux5e8f}

\begin{enumerate}
\def\labelenumi{\arabic{enumi}.}
\tightlist
\item
  \textbf{3D Location因果决定标签 + None冗余Structure编码}（CNT
  手性、缺陷空间分辨）→ \(\delta_s^{\text{PE}}\) 大且正。这Yes Situs
  的设计Objective场景。
\item
  \textbf{3D Location因果决定标签 +
  有冗余但有额外Structure}（对接评分、Materials缺陷）→ \(\delta_s^{\text{PE}}\)
  中且正，但DimensionSelection和正则化关键。
\item
  \textbf{1D 序列Location弱相关标签 + NoneStructure编码}（easy蛋白质功能预测，None
  pLM）→ \(\delta_s^{\text{PE}}\) 小且正，微弱收益。
\end{enumerate}

\subsection{它何时None用（加了等于白加）}\label{ux5b83ux4f55ux65f6ux65e0ux7528ux52a0ux4e86ux7b49ux4e8eux767dux52a0}

\begin{enumerate}
\def\labelenumi{\arabic{enumi}.}
\setcounter{enumi}{3}
\tightlist
\item
  \textbf{Location不携带给定Status的额外标签信息}（IDR、纯组成分类、体相性质）→
  \(\delta_s^{\text{PE}} \approx 0\)，纯Dimension浪费。
\item
  \textbf{已有强力Structure编码器}（GNN 在晶体上、pLM 在蛋白上）→
  共线性，\(\delta_s^{\text{PE}} \approx 0\) 但有微小负效应。
\end{enumerate}

\subsection{它何时有害（加了更差）}\label{ux5b83ux4f55ux65f6ux6709ux5bb3ux52a0ux4e86ux66f4ux5dee}

\begin{enumerate}
\def\labelenumi{\arabic{enumi}.}
\setcounter{enumi}{5}
\tightlist
\item
  \textbf{伪Location或随机Location +
  有限Sample}（NoneStructureData的虚拟索引、随机构型）→ \(\delta_s^{\text{PE}}\)
  Training \textgreater{} 0 但Test \textless{} 0，最危险的虚假相关。
\item
  \textbf{标签具有平移/旋转对称性但任务不尊重该对称性}（体相性质 + Situs
  编码打破对称性）→ 引入与Physics矛盾的归纳偏差。
\item
  \textbf{未修正的Computation artifact
  被编码为''Location''}（混合不同超胞尺寸的未修正 DFT Data）→ 学到Computation
  artifact 而非Physics。
\end{enumerate}

\subsection{Theory洞察与工程现实的鸿沟}\label{ux7406ux8bbaux6d1eux5bdfux4e0eux5de5ux7a0bux73b0ux5b9eux7684ux9e3fux6c9f}

定理 2.2.1 给出的充分Conditions \(I(Y; P \mid S) > 0\)
Yes\textbf{必要且精确的}，但在实践中： - 估计 \(I(Y; P \mid S)\)
本身需要大量Sample（KSG 估计器的ConvergenceDimension灾难，定理 3.3.1) - 即使
\(I(Y; P \mid S) > 0\)，定理 2.3.1 的Upper Bound可能因编码不完美度
\(\varepsilon_{\text{PE}}\) 而极度松弛 - Theorem 3' 揭示的''Learning型 PPE
可能区分Noise和DifficultySample''在实践中\textbf{通常不发生}------因为Physics学常用的
PE Yes固定的，而Learning型 PE 的信号太弱None法可靠破坏不可区分性

\begin{center}\rule{0.5\linewidth}{0.5pt}\end{center}

\emph{AnalysisDate：2026-06-29}\\
\emph{PhysicsData引用：Bartlett (2002), Creighton (1993), Dunker (2005),
Freysoldt (2014), Makov \& Payne (1995), Pauling (1951), Stampfl \& Van
de Walle (2002), Van de Walle \& Neugebauer (2004), Yan (2016), 等}\\
\emph{Methodology声明：所有 \(\delta_s^{\text{PE}}\)
的估计来自Physics推演而非数值Experiment------需要 DFT
Computation和蛋白质StructureAnalysis来验证/校准这些估计。}

\end{document}
